\documentclass[../../../include/open-logic-section]{subfiles}

\begin{document}

\olfileid[hi]{sth}{replacement}{strength}
\olsection{प्रतिस्थापन की शक्ति}

प्रतिस्थापन की शक्ति के विषय में एक सरल प्रेक्षण से आरंभ करें: जब तक हम
$\Z$ से आगे नहीं जाते, $\omega + \omega$ से बड़े या उसके बराबर किसी वॉन
नोयमान क्रमसूचक का अस्तित्व सिद्ध नहीं कर सकते।

कारण की रूपरेखा यह है।~$\ZF$ में काम करते हुए समुच्चय
$V_{\omega+\omega}$ पर विचार करें। यह समुच्चय~$\Z$ के किसी \emph{प्रतिरूप}
के प्रांत का कार्य करता है। इसे देखने के लिए हम सूत्र के
\emph{सापेक्षीकरण} का कुछ संकेतन प्रस्तुत करते हैं:
\begin{defn}\ollabel{formularelativization}
	किसी भी समुच्चय $M$ और किसी भी सूत्र $\phi$ के लिए $\phi^M$ को वह सूत्र
	मानें जो $\phi$ के सभी परिमाणकों को $M$ तक सीमित करने से मिलता है। अर्थात
	``$\lexists[x]$'' को ``$(\lexists[x \in M])$'' से और
	``$\lforall[x]$'' को ``$(\lforall[x \in M])$'' से बदलें।
\end{defn}\noindent
दिखाया जा सकता है कि~$\Z$ के प्रत्येक अभिगृहीत $\phi$ के लिए
$\ZF \vdash \phi^{V_{\omega+\omega}}$। किंतु
\olref[spine][rank]{ordsetrankalpha} से $\omega+\omega$,
$V_{\omega+\omega}$ में \emph{नहीं} है। अतः $\omega+\omega$ का अनस्तित्व
$\Z$ के साथ संगत है।

इसीलिए हमने \olref[ordinals][replacement]{sec} में कहा था कि प्रतिस्थापन के
बिना \olref[ordinals][ordtype]{thmOrdinalRepresentation} सिद्ध नहीं किया जा
सकता। क्योंकि~$\Z$ में ऐसी स्पष्ट सुव्यवस्था परिभाषित करना सरल है जिसका सहज
रूप से क्रम-प्रकार $\omega+\omega$ \emph{होना चाहिए}। वास्तव में
\olref[ordinals][idea]{sec} में हमने प्राकृतिक संख्याओं पर इस क्रम का
अनौपचारिक उदाहरण दिया था:
\begin{align*}
	n \lessdot m \text{ तभी और केवल तभी जब }&\text{या तो }n < m\text{ और }m-n\text{ सम है,}\\
	& \text{अथवा $n$ सम और $m$ विषम है।}
\end{align*}
किंतु यदि $\omega+\omega$ का अस्तित्व नहीं, तो यह सुव्यवस्था किसी क्रमसूचक
के समाकृतिक नहीं। अतः $\Z$,
\olref[ordinals][ordtype]{thmOrdinalRepresentation} को सिद्ध \emph{नहीं}
करता।

बात को उलटें: प्रतिस्थापन हमें $\omega+\omega$ का अस्तित्व सिद्ध करने देता
है और इसलिए $V_{\omega+\omega}$ का अस्तित्व भी सिद्ध करने देना चाहिए। और
इतना ही नहीं। हम जो \emph{कोई भी} सुव्यवस्था परिभाषित कर सकें, उसके लिए
\olref[ordinals][ordtype]{thmOrdinalRepresentation} बताता है कि कोई
$\alpha$ उस सुव्यवस्था के समाकृतिक है और इसलिए $V_\alpha$ का अस्तित्व है।
इस सीधे अर्थ में प्रतिस्थापन गारंटी देता है कि समुच्चयों का पदानुक्रम
\emph{बहुत ऊँचा} होना चाहिए।

अगले कुछ अनुभागों में, और फिर \olref[card-arithmetic][fix]{sec} में, हमें
अधिक स्पष्ट आभास होगा कि प्रतिस्थापन पदानुक्रम को ठीक \emph{कितना} ऊँचा
होने को बाध्य करता है। अभी सरल बात यह है कि प्रतिस्थापन को सचमुच औचित्य की
\emph{आवश्यकता है}!{}

\end{document}
